\documentclass[a4paper,12pt]{article}

% Paquetes necesarios
\usepackage[utf8]{inputenc}   % Codificación UTF-8
\usepackage{fontspec} % Paquete para manejar fuentes específicas
\usepackage[spanish]{babel}
\usepackage[style=apa, backend=biber]{biblatex}
\usepackage{geometry}         % Configuración de márgenes
\usepackage{graphicx}
\usepackage{setspace}         % Configuración de interlineado
\usepackage{fancyhdr}         % Configuración de encabezados y pie de página

\setmainfont{Arial}   % Configura Arial como la fuente principal
\geometry{
    top=2.5cm,
    bottom=2.5cm,
    left=4cm,
    right=2.5cm
}
\pagestyle{fancy}
% Configuración del interlineado y espacio entre párrafos
\setlength{\parskip}{2em}  % Espacio doble entre párrafos
\setlength{\parindent}{0em} % Sin sangría en párrafos
\onehalfspacing            % Interlineado de 1.5
% Configuración del pie de página
\fancyhf{}
\fancyfoot[R]{\thepage}    % Número de página en la esquina inferior derecha
\fancypagestyle{plain}{
    \fancyhf{}
    \fancyfoot[R]{\thepage}
}


\addbibresource{bibliography.bib} % Archivo .bib con las referencias



\begin{document}

% Incluir la portada
% Archivo: portada.tex
\begin{titlepage}
    \begin{center}
        {\fontsize{15pt}{10pt}\selectfont\textbf{UNIVERSIDAD AUTÓNOMA "GABRIEL RENE MORENO"}} \\
        {\fontsize{15pt}{10pt}\selectfont\textbf{FACULTAD DE INGENIERÍA EN CIENCIAS DE LA COMPUTACIÓN Y TELECOMUNICACIONES}} \\
        \vspace{1cm}
        % Logo
        \includegraphics[width=0.8\textwidth]{images/logo_soe.png} \\
        
        % Título de la institución
        {\fontsize{15pt}{10pt}\selectfont\textbf{MAESTRÍA EN CIENCIA DE DATOS E INTELIGENCIA ARTIFICIAL}} \\

        % Título del trabajo
        \vspace{1cm}
        {\fontsize{15pt}{5pt}\selectfont\textbf{Implementar un Chatbot Semántico Multiperfil para Acceso Inteligente a Documentación Técnica y de Negocio en Plataformas Confluence del Sector Aeronáutico en Airnguru S.A. para la Gestión 2025}} \\
        \vspace{1cm}
        
        \vspace{2cm}
        \begin{flushright}
        {\fontsize{14pt}{22pt}\selectfont\textbf{Autor}} \\
        {\fontsize{14pt}{22pt}\selectfont{Ing. Darlyn Bravo Peña}} \\
        \end{flushright}
        
        
        
        % Ciudad y fecha
        \vfill
        \textbf{Santa Cruz, Bolivia} \\
    \end{center}
    \end{titlepage}
     % Reemplaza con % Archivo: portada.tex
\begin{titlepage}
    \begin{center}
        {\fontsize{15pt}{10pt}\selectfont\textbf{UNIVERSIDAD AUTÓNOMA "GABRIEL RENE MORENO"}} \\
        {\fontsize{15pt}{10pt}\selectfont\textbf{FACULTAD DE INGENIERÍA EN CIENCIAS DE LA COMPUTACIÓN Y TELECOMUNICACIONES}} \\
        \vspace{1cm}
        % Logo
        \includegraphics[width=0.8\textwidth]{images/logo_soe.png} \\
        
        % Título de la institución
        {\fontsize{15pt}{10pt}\selectfont\textbf{MAESTRÍA EN CIENCIA DE DATOS E INTELIGENCIA ARTIFICIAL}} \\

        % Título del trabajo
        \vspace{1cm}
        {\fontsize{15pt}{5pt}\selectfont\textbf{Implementar un Chatbot Semántico Multiperfil para Acceso Inteligente a Documentación Técnica y de Negocio en Plataformas Confluence del Sector Aeronáutico en Airnguru S.A. para la Gestión 2025}} \\
        \vspace{1cm}
        
        \vspace{2cm}
        \begin{flushright}
        {\fontsize{14pt}{22pt}\selectfont\textbf{Autor}} \\
        {\fontsize{14pt}{22pt}\selectfont{Ing. Darlyn Bravo Peña}} \\
        \end{flushright}
        
        
        
        % Ciudad y fecha
        \vfill
        \textbf{Santa Cruz, Bolivia} \\
    \end{center}
    \end{titlepage}
     si prefieres que se trate como una sección independiente

% Table of contents
\newpage
\tableofcontents
\listoffigures

\begin{center}
    {\LARGE \textbf{INTRODUCCIÓN}} % Título centrado y en negrita
\end{center}

La estimación de costos en el desarrollo de software es un desafío inherente a la naturaleza intangible de este producto. En un mundo donde la economía digital impulsa la innovación y la competitividad, la incapacidad de calcular con precisión los costos de desarrollo puede llevar a pérdidas financieras, proyectos fallidos y relaciones comerciales insatisfactorias. A nivel global, empresas de todos los tamaños enfrentan la necesidad de establecer modelos más eficientes y estandarizados para valorar su trabajo, especialmente en un mercado donde la externalización y la globalización son cada vez más comunes. Este tema adquiere mayor relevancia 
con el auge de las metodologías ágiles y la inteligencia artificial, que prometen pero también complican la predicción de costos.

En Bolivia, donde el sector tecnológico está en crecimiento, el problema de la estimación de costos en el desarrollo de software adquiere un matiz especial. La falta de modelos estandarizados y herramientas locales para valorar los servicios de desarrollo de software lleva a una alta dependencia de las negociaciones subjetivas, lo que puede perjudicar tanto a los desarrolladores como a los clientes. Esto es especialmente importante en un contexto donde muchas empresas están comenzando a digitalizarse y dependen del software para competir en un mercado nacional e internacional. Abordar este desafío no solo contribuirá al fortalecimiento del sector tecnológico en Bolivia, sino que también permitirá que los desarrolladores y empresas logren acuerdos más justos y sostenibles, impulsando así el crecimiento de la economía digital en el país.

El tema está ubicado en el area 
Ciencias Generales de la Computación Aplicadas y la 
línea Ingeniería de software con un eje de 
Aspectos económicos y de negocio del proceso de desarrollo del software porque la estimación de costos afecta la 
rentabiliad de los proyectos del software, tanto para los desarrolladores
como para los clientes.
\section{Antecedentes}
La estimación precisa de costos en el desarrollo de software es un desafío constante en la ingeniería de software, especialmente en entornos ágiles. Técnicas como los \textit{Puntos de Historia} y el \textit{Planning Poker} han sido desarrolladas para abordar este problema, permitiendo a los equipos estimar el esfuerzo necesario para completar historias de usuario de manera colaborativa 
y basada en consenso \parencite{asana_planning_poker}
Los \textit{Puntos de Historia} son una unidad de medida abstracta que refleja el esfuerzo relativo, la complejidad y el riesgo asociados a una historia de usuario. Esta técnica permite a los equipos evaluar tareas en función de su tamaño relativo, en lugar de estimar en unidades de tiempo, lo que facilita 
la planificación y priorización en entornos ágiles \parencite{asana_story_points}.

Por otro lado, el \textit{Planning Poker} es una técnica de estimación que combina la opinión de todos los miembros del equipo para alcanzar una estimación consensuada. Durante una sesión de \textit{Planning Poker}, cada miembro selecciona de manera independiente una carta con el valor que considera adecuado para una tarea específica. Posteriormente, las discrepancias se discuten hasta alcanzar un consenso, promoviendo 
la colaboración y la precisión en las estimaciones \parencite{asana_planning_poker}.

A pesar de su efectividad en proyectos ágiles, estas técnicas enfrentan desafíos, como la subjetividad inherente al proceso y la falta de estandarización en algunos equipos. Sin embargo, su uso ha demostrado ser fundamental para mejorar la precisión de las estimaciones y la planificación 
en proyectos de software de diversa escala \parencite{asana_planning_poker}.


\printbibliography

\end{document}
